% ======================================================================
\section{Setup: soft gluon wave function and the evolution operator}\label{sec:setup}
% ======================================================================

In this section we summarise the framework of~\cite{Radhakrishnan:2026soft}, on which the present calculation is built, and collect the coefficients of the evolution operator $\Omega$ that enter the single inclusive cross section. The coefficients of order $g$ and $g^2$ were derived in~\cite{Radhakrishnan:2026soft}; the coefficient of order $g^3$ that creates one gluon, $\mf A^{(3)}_2$, is obtained here. The derivations of the light-cone wave functions from which the coefficients are obtained are summarised in Appendix~\ref{app:wavefunctions}.

\subsection{Light-cone conventions}\label{sec:conventions}

We use light-cone coordinates $x^\pm=(x^0\pm x^3)/\sqrt2$ and denote transverse vectors by $\v x=(x^1,x^2)$; where no confusion is possible the transverse vectors are written without boldface, and $\int_{x}\equiv\int d^2x$. The projectile moves in the $+$ direction, $x^+$ plays the role of time, and $p^+>0$ is the longitudinal momentum of a physical particle. The light-cone energy of an on-shell gluon is $p^-=\v p^2/2p^+$, and we write
\begin{align}
&E_g(p)=\frac{\v p^2}{2p^+}\,,\qquad E_{gg}(p,q)=\frac{\v p^2}{2p^+}+\frac{\v q^2}{2q^+}\,,\notag\\ &d^3p\equiv dp^+d^2p\,,\qquad \delta^{(3)}(p-q)\equiv\delta(p^+-q^+)\,\delta^{(2)}(\v p-\v q)\,.
\end{align}
We work in the light-cone gauge $A^+=0$, in which the transverse components $A_i^a$, $i=1,2$, are the dynamical fields. The gluon creation and annihilation operators satisfy
\begin{align}\label{eq:aadagger}
\big[a^a_i(k),a^{\dagger b}_j(p)\big]=(2\pi)^3\,\delta^{ab}\,\delta_{ij}\,\delta^{(3)}(k-p)\,.
\end{align}
For the Fourier transforms in the transverse plane we use the conventions
\begin{align}\label{eq:FTconventions}
&a(k^+,\v k)=\int d^2z\,e^{-i\v k\cdot\v z}\,a(k^+,\v z)\,,\qquad
a^\dagger(k^+,\v k)=\int d^2z\,e^{i\v k\cdot\v z}\,a^\dagger(k^+,\v z)\,,\notag\\
&f(\v k)=\int d^2x\,e^{i\v k\cdot\v x}\,f(\v x)\,,
\end{align}
so that $\big[a^a_i(k^+,\v z),a^{\dagger b}_j(p^+,\v z')\big]=2\pi\,\delta^{ab}\delta_{ij}\,\delta(k^+-p^+)\,\delta^{(2)}(\v z-\v z')$. One- and two-gluon states are normalised as
\begin{align}\label{eq:states}
\ket{g^a_i(k^+,\v z)}=\frac{a^{\dagger a}_i(k^+,\v z)}{(2\pi)^{1/2}}\ket{0}\,,\qquad
\ket{g^a_i(k^+,\v z)\,g^b_j(p^+,\v z')}=\frac{a^{\dagger a}_i(k^+,\v z)\,a^{\dagger b}_j(p^+,\v z')}{2\pi}\ket{0}\,,
\end{align}
with $\braket{g^b_j(p^+,\v z')|g^a_i(k^+,\v z)}=\delta^{ab}\delta_{ij}\,\delta(k^+-p^+)\,\delta^{(2)}(\v z-\v z')$; in momentum space $\ket{g^a_i(k)}=(2\pi)^{-3/2}a^{\dagger a}_i(k)\ket0$ (Appendix~\ref{app:states}).

Throughout the paper we use the Weizs\"acker--Williams field of a point charge and the transverse kernels built from it,
\begin{align}\label{eq:kernels}
\Kv^i(\v x)\equiv\frac{x^i}{\v x^2}\,,\qquad
K(\v x,\v y,\v z)\equiv\Kv(\v x-\v z)\cdot\Kv(\v y-\v z)=\frac{(\v x-\v z)\cdot(\v y-\v z)}{(\v x-\v z)^2(\v y-\v z)^2}\,,
\end{align}
so that $K(\v x,\v x,\v z)=1/(\v x-\v z)^2$ and $K(\v x,\v x,\v z)-2K(\v x,\v y,\v z)+K(\v y,\v y,\v z)=(\v x-\v y)^2/[(\v x-\v z)^2(\v y-\v z)^2]$ is the dipole kernel.

\subsection{Soft and valence modes; the eikonal Hamiltonian}\label{sec:hamiltonian}

Following the CGC logic~\cite{McLerran:1993ni,McLerran:1993ka,Jalilian-Marian:1997jx,Kovner:2005en}, the gluon modes of the projectile are divided according to their longitudinal momentum. The valence modes have $k^+>\vee$, the soft modes have
\begin{align}
\Lambda<k^+<\vee\,,\qquad \vee\equiv\Lambda\,e^{\delta Y}\,,
\end{align}
and the modes with $k^+<\Lambda$ do not belong to the projectile wave function; their effect is part of the target. The rapidity interval $\delta Y=\ln(\vee/\Lambda)$ is the interval over which the soft modes are integrated; the notation is that of~\cite{Lublinsky:2016meo,Radhakrishnan:2026soft}. Since the valence modes are much faster than the soft ones, they act on the soft modes only through their colour charge density
\begin{align}\label{eq:rho}
\rho^a(\v x)\,,\qquad \big[\rho^a(\v x),\rho^b(\v y)\big]=i f^{abc}\,\rho^c(\v x)\,\delta^{(2)}(\v x-\v y)\,,\qquad
\big[\rho^a(\v k),\rho^b(\v p)\big]=i f^{abc}\,\rho^c(\v k+\v p)\,,
\end{align}
which is an operator on the valence Hilbert space; its microscopic definition in terms of the valence gluon operators is given in~\cite{Radhakrishnan:2026soft}. We keep the charges non-commuting throughout. This is necessary: products of $\rho$'s appear in the cross section in a definite order, determined by the time ordering of the light-cone perturbation theory, and the commutators~\eqref{eq:rho} reduce the number of $\rho$'s by one. These reductions produce some of the logarithmic terms of Section~\ref{sec:rapidity}.

After the decomposition of the gauge field into valence and soft components, all terms suppressed by powers of $k^+_{\rm soft}/k^+_{\rm valence}$ are dropped (eikonal approximation). The light-cone Yang--Mills Hamiltonian ($N_f=0$) then takes the form~\cite{Lublinsky:2016meo,Radhakrishnan:2026soft}
\begin{align}\label{eq:Hdecomp}
&H=H_0+H_g+H_{ggg}+H_{gg\text{-inst}}+H_{gggg}+\dots,\notag\\
&H_0=\int_\Lambda^\vee\frac{dk^+}{2\pi}\int\frac{d^2k}{(2\pi)^2}\,\frac{\v k^2}{2k^+}\,a^{\dagger a}_i(k^+,\v k)\,a^a_i(k^+,\v k)\,,
\end{align}
where the dots stand for terms which do not contribute at the order considered. The coupling of the soft gluons to the valence charges is
\begin{align}\label{eq:Hg}
H_g=\int_\Lambda^\vee\frac{dk^+}{2\pi}\int\frac{d^2k}{(2\pi)^2}\,\frac{g\,k^i}{\sqrt2\,|k^+|^{3/2}}\Big[a^{\dagger a}_i(k^+,\v k)\,\rho^a(-\v k)+a^a_i(k^+,\v k)\,\rho^a(\v k)\Big]\,.
\end{align}
It describes the emission and absorption of a soft gluon by the valence charges. $H_{ggg}$ is the triple-gluon vertex among the soft modes, $H_{gg\text{-inst}}$ the instantaneous (Coulomb-like) interaction between two soft gluons and the valence charge density, and $H_{gggg}$ the four-gluon interaction of the soft modes. Their explicit forms and the matrix elements used in this paper are listed in Appendix~\ref{app:hamiltonian}.

\subsection{Light-cone perturbation theory with operator-valued matrix elements}\label{sec:lcpt}

The soft gluon wave function of the projectile is obtained in time-independent (old-fashioned) perturbation theory in $H-H_0$ on the light cone~\cite{Lepage:1980fj,Brodsky:1997de}. Since the matrix elements of $H_g$ and $H_{gg\text{-inst}}$ are operators on the valence space, the perturbative series is written with the matrix elements in the order of the light-cone time ordering: the vertex acting later stands to the left. With $E_0=0$ for the soft vacuum, the first terms of the expansion of the dressed state in terms of free Fock states are
\begin{align}\label{eq:lcpt}
\ket{\Psi}=\ &\ket0-\sum_n\ket n\frac{V_{n0}}{E_n}+\sum_{n,m}\ket n\frac{V_{nm}V_{m0}}{E_nE_m}
-\sum_{n,m,l}\ket n\frac{V_{nm}V_{ml}V_{l0}}{E_nE_mE_l}
+\sum_{n,m}\ket n\frac{V_{n0}\,V_{0m}V_{m0}}{E_n^2E_m}
\notag\\
&-\frac12\sum_m\frac{V_{0m}V_{m0}}{E_m^2}\ket0
+\frac12\sum_{n,m}\ket n\frac{V_{n0}\,V_{0m}V_{m0}}{E_nE_m^2}+\dots,
\end{align}
where $V_{nm}=\braket{n|H-H_0|m}$ and the sums run over Fock states of the soft sector. The first line is the Rayleigh--Schr\"odinger series with intermediate normalisation; the second line implements the canonical normalisation of the state and is needed at orders $g^2$ (vacuum term) and $g^3$ (one-gluon term). In the last term of each line the single-emission matrix element $V_{n0}$ stands to the left of the norm factor $V_{0m}V_{m0}$; this ordering follows from the perturbative construction of the unitary operator that diagonalises $H$ with operator-valued couplings, and we have verified it against an explicit construction of $\Omega$ with finite-dimensional matrices representing the charges. Energy denominators that vanish inside an integration region are understood as principal values (Appendix~\ref{app:distributions}).

\subsection{The evolution operator}\label{sec:omega}

The dressed (physical) states of the projectile are obtained from the free Fock states by a unitary operator $\Omega$, $\Omega^\dagger\Omega=1$. Since the soft modes are treated in the eikonal approximation, $\Omega$ can be written as a normal-ordered series in the soft gluon operators, with coefficients that are functionals of $\rho$~\cite{Radhakrishnan:2026soft}:
\begin{align}\label{eq:OmegaExp}
\Omega = \ &\mf{N}
+\int\frac{d^3p}{(2\pi)^3}\Big[\mf{A}^b_{1j}(p)\,a^b_j(p)+\mf{A}^b_{2j}(p)\,a^{\dagger b}_j(p)\Big]
\notag\\&
+\int\frac{d^3p}{(2\pi)^3}\frac{d^3q}{(2\pi)^3}\Big[\mf{B}^{cb}_{1kj}(p,q)\,a^b_j(p)a^c_k(q)+\mf{B}^{cb}_{2kj}(p,q)\,a^{\dagger b}_j(p)a^{\dagger c}_k(q)
\notag\\&\hspace{3.4cm}+\mf{B}^{cb}_{3kj}(p,q)\,a^{\dagger c}_k(q)a^b_j(p)\Big]
\notag\\&
+\int\frac{d^3p}{(2\pi)^3}\frac{d^3q}{(2\pi)^3}\Big[\mf{C}^{dcb}_{1lkj}(p,q)\,a^{\dagger d}_l(q)a^{\dagger c}_k(p-q)a^b_j(p)
\notag\\&\hspace{3.4cm}+\mf{C}^{dcb}_{2lkj}(p,q)\,a^{\dagger b}_j(p)a^c_k(p-q)a^d_l(q)\Big]+\dots
\end{align}
The dots denote terms with more gluon operators, which do not contribute to the single inclusive cross section at $O(g^4)$. The coefficients have definite orders in $g$: $\mf N=1+\mf N_2$ with $\mf N_2=O(g^2)$; $\mf A_{1,2}=\mf A^{(1)}_{1,2}+\mf A^{(3)}_{1,2}$ with $\mf A^{(n)}=O(g^n)$; $\mf B_{1,2,3}=O(g^2)$; $\mf C_{1,2}=O(g)$. The coefficients that multiply only creation operators, $\mf A_2$ and $\mf B_2$, together with $\mf N$, are fixed by the action of $\Omega$ on the soft vacuum, i.e.\ by the light-cone wave function of the projectile; $\mf B_3$ and $\mf C_1$ are fixed by the action of $\Omega$ on a one-gluon state. The remaining coefficients follow from unitarity, imposed order by order in $g$ (Appendix~\ref{app:unitarity}):
\begin{align}\label{eq:unitarity_short}
&\mf A^{(1)\dagger b}_{1j}(p)=-\mf A^{(1)b}_{2j}(p)\,,\qquad
\mf C^{\dagger dcb}_{2lkj}(p,q)=-\mf C^{dcb}_{1lkj}(p,q)\,,\notag\\
&\mf N=1-\frac12\int\frac{d^3p}{(2\pi)^3}\,\mf A^{(1)b}_{1j}(p)\,\mf A^{(1)\dagger b}_{1j}(p)\,,
\end{align}
together with the relations for $\mf A^{(3)}_1$ and $\mf B_1$ given in Eqs.~\eqref{eq:A31unit} and~\eqref{eq:B1unit}. In the relation for $\mf A^{(3)}_1$ the operator character of the coefficients matters: $\mf N_2$ does not commute with $\mf A^{(1)}$, since $[\rho^a(\v x)\rho^a(\v y),\rho^b(\v w)]\neq0$ for $\v x\ne\v y$.

In transverse coordinate space, after the unitarity relations~\eqref{eq:unitarity_short} are used, the part of $\Omega$ that contributes is
\begin{align}\label{eq:OmegaX}
\Omega = \ &\mf N
+\int\frac{dp^+}{2\pi}\int_x\Big[\mf A^{(1)b}_{1j}(p^+,\v x)\,a^b_j(p^+,\v x)-\mf A^{\dagger(1)b}_{1j}(p^+,\v x)\,a^{\dagger b}_j(p^+,\v x)
\notag\\&\hspace{2.6cm}
+\mf A^{(3)b}_{1j}(p^+,\v x)\,a^b_j(p^+,\v x)+\mf A^{(3)b}_{2j}(p^+,-\v x)\,a^{\dagger b}_j(p^+,\v x)\Big]
\notag\\&
+\int\frac{dp^+}{2\pi}\frac{dq^+}{2\pi}\int_{x,y}\Big[\mf B^{cb}_{1kj}(p^+,\v x;q^+,\v y)\,a^b_j(p^+,\v x)a^c_k(q^+,\v y)
\notag\\&\hspace{3.2cm}
+\mf B^{cb}_{2kj}(p^+,-\v x;q^+,-\v y)\,a^{\dagger b}_j(p^+,\v x)a^{\dagger c}_k(q^+,\v y)
\notag\\&\hspace{3.2cm}
+\mf B^{cb}_{3kj}(p^+,\v x;q^+,-\v y)\,a^{\dagger c}_k(q^+,\v y)a^b_j(p^+,\v x)\Big]
\notag\\&
+\int\frac{dp^+}{2\pi}\frac{dq^+}{2\pi}\int_{x,y,z}\Big[\mf C^{dcb}_{1lkj}(p^+,\v z-\v y;q^+,\v y-\v x)
\notag\\&\hspace{3.2cm}\times a^{\dagger d}_l(q^+,\v x)a^{\dagger c}_k(p^+-q^+,\v y)a^b_j(p^+,\v z)
\notag\\&\hspace{3.2cm}
-\mf C^{\dagger dcb}_{1lkj}(p^+,\v z-\v y;q^+,\v y-\v x)
\notag\\&\hspace{3.2cm}\times a^{\dagger b}_j(p^+,\v z)a^c_k(p^+-q^+,\v y)a^d_l(q^+,\v x)\Big]+\dots
\end{align}
The transverse arguments of the coefficients are fixed by the Fourier conventions~\eqref{eq:FTconventions}.

\subsection{The coefficients}\label{sec:coefficients}

We now list the coefficients in transverse coordinate space. They are obtained by comparing $\Omega\ket0$ and $\Omega\ket{g^a_i(k^+,\v z)}$ with the LCPT wave functions of Appendix~\ref{app:wavefunctions}, sector by sector in Fock space; the factors of $2\pi$ come from the normalisation~\eqref{eq:states}. The momentum-space expressions and the intermediate steps are given in Appendix~\ref{app:wavefunctions} (for the coefficients of order $g$ and $g^2$ also in~\cite{Radhakrishnan:2026soft}).

\paragraph{Normalisation.} The coefficient of the identity is fixed by the norm of the LO wave function,
\begin{align}\label{eq:N2}
\mf N=1+\mf N_2\,,\qquad
\mf N_2&=-\frac{g^2}{8\pi^3}\,\ln\frac{\vee}{\Lambda}\int_{x,y,z}K(\v x,\v y,\v z)\,\rho^a(\v x)\rho^a(\v y)
\notag\\&=-\frac{g^2}{8\pi^3}\,\ln\frac{\vee}{\Lambda}\int d^2p\,\frac{\rho^a(-\v p)\rho^a(\v p)}{\v p^2}\,.
\end{align}
$\mf N_2$ is the (virtual) probability of not emitting a soft gluon anywhere in the interval $\Lambda<p^+<\vee$, and it carries the full logarithm $\ln(\vee/\Lambda)$.

\paragraph{Single-gluon emission at $O(g)$.} The Weizs\"acker--Williams amplitude for the emission of a soft gluon at $\v z$ is
\begin{align}\label{eq:A1}
\mf A^{(1)a}_{2i}(k^+,-\v z)=-\mf A^{(1)\dagger a}_{1i}(k^+,\v z)=\frac{ig}{\sqrt2\,\pi\sqrt{k^+}}\int_x\Kv^i(\v x-\v z)\,\rho^a(\v x)\,.
\end{align}
It corresponds to the LO wave function $\ket{\Psi^{\rm LO}}=\frac{ig}{2\pi^{3/2}}\int_\Lambda^\vee\frac{dk^+}{\sqrt{k^+}}\int_{x,z}\Kv^i(\v x-\v z)\,\rho^a(\v x)\ket{g^a_i(k^+,\v z)}$, Eq.~\eqref{eq:PsiLO}. In what follows we use the shorthand $\mf A\equiv\mf A^{(1)}_2$ when no confusion can arise, and a bar denotes the rotation of all charges by the target, $\bar{\mf A}^a(\v x)=\mf A^a(\v x)\big|_{\rho(\v u)\to U(\v u)\rho(\v u)}$, see Section~\ref{sec:observable}.

\paragraph{Single-gluon emission at $O(g^3)$.} The coefficient $\mf A^{(3)}_2$ contains terms with one, two and three charges. In momentum space the one-$\rho$ part reads, for $\Lambda\to0$ wherever this is possible,
\begin{align}\label{eq:A3onerho}
\Big[\mf A^{(3)a}_{2i}(k^+,\v k)\Big]_{\rho}=\ &\frac{g^3N_c\,\rho^a(-\v k)\,k^i}{8\sqrt{2k^+}\,\pi^2\,\v k^2}\Bigg\{\Big[-\frac2\epsilon+\ln\frac{\v k^2}{\mu^2_{\MS}}\Big]
\notag\\&\hspace{1.5cm}\times\Big[\frac{11}{3}-2\ln\frac{k^+}{\Lambda}+\ln\Big(1-\frac{k^+}{\vee}\Big)\Big]
\notag\\&\hspace{1.5cm}
+\frac{\pi^2}{3}-\frac{67}{9}+\ln^2\Big(1-\frac{k^+}{\vee}\Big)+\ln\Big(1-\frac{k^+}{\vee}\Big)\ln\frac{\vee}{k^+}\Bigg\},
\end{align}
where the transverse integrals are performed in $d=2-\epsilon$ dimensions in the $\MS$ scheme. It collects the one-loop self-energy of the emitted gluon, the vertex corrections in which the emitted gluon splits and recombines, and the one-$\rho$ (commutator) parts of the diagrams with two valence vertices (Figures~\ref{fig:vac_g3_onerho} and~\ref{fig:vac_g3_tworho}). The pole multiplies the combination $\frac{11}{3}-2\ln(k^+/\Lambda)$, whose first term is $\beta_0/N_c$ for pure glue, $\beta_0=\frac{11}{3}N_c$. In coordinate space,
\begin{align}\label{eq:A3onerhoX}
\Big[\mf A^{(3)a}_{2i}(k^+,-\v z)\Big]_{\rho}=\ &-\int_x\frac{ig^3N_c\,\rho^a(\v x)\,\Kv^i(\v x-\v z)}{16\sqrt2\,\pi^3\sqrt{k^+}}\Bigg\{\Big[-\frac2\epsilon-2\gamma_E-\ln\frac{(\v x-\v z)^2\mu^2_{\MS}}{4}\Big]
\notag\\&\hspace{1.5cm}\times\Big[\frac{11}{3}-2\ln\frac{k^+}{\Lambda}+\ln\Big(1-\frac{k^+}{\vee}\Big)\Big]
\notag\\&\hspace{1.5cm}
-\frac{67}{9}+\frac{\pi^2}{3}+\ln^2\Big(1-\frac{k^+}{\vee}\Big)+\ln\Big(1-\frac{k^+}{\vee}\Big)\ln\frac{\vee}{k^+}\Bigg\}.
\end{align}
The two-$\rho$ part is antisymmetric in colour, $\propto f^{abc}\{\rho^b,\rho^c\}$, and is conveniently written in a mixed representation,
\begin{align}\label{eq:A3tworho}
\Big[\mf A^{(3)a}_{2i}(k^+,-\v z)\Big]_{\rho\rho}=\ &\int d^2p\,d^2k\int_{x,y}e^{i\v k\cdot(\v z-\v y)}e^{i\v p\cdot(\v y-\v x)}\,\frac{ig^3f^{abc}\{\rho^b(\v x),\rho^c(\v y)\}}{32\sqrt{2k^+}\,\pi^5\,\v p^2\,\v k^2}
\notag\\&\times\Bigg[\ln\frac{\Lambda}{k^+}\,\frac{\v k^2p^i-(\v k\cdot\v p)\,k^i}{(\v k-\v p)^2}
-2k^i\ln\frac{\vee}{\Lambda}
-2k^i\ln\frac{k^+\v p^2}{\vee\,\v k^2+k^+\v p^2}
\notag\\&\quad
+\frac{(2\v p^2-2\v k\cdot\v p+\v k^2)\,k^i+\v k^2p^i}{(\v k-\v p)^2}\ln\Big[1+\Big(\frac{\vee}{k^+}-1\Big)\frac{\v k^2}{(\v k-\v p)^2}\Big]
\notag\\&
-\frac{\v p^2k^i+\v k^2p^i}{(\v k-\v p)^2}\ln\frac{\vee-k^+}{\Lambda}-k^i\ln\frac{\vee}{k^+}\Bigg],
\end{align}
where, inside the integral, $\v k$ and $\v p$ are integration variables. The three-$\rho$ part originates from the third-order terms of Eq.~\eqref{eq:lcpt} in which a gluon is emitted and reabsorbed by the valence charges before or after the emission of the observed gluon (Figure~\ref{fig:vac_g3_threerho}). The terms proportional to $1/\Lambda$ and $1/\vee$ cancel exactly, and the remainder is
\begin{align}\label{eq:A3threerho}
\Big[\mf A^{(3)a}_{2i}(k^+,-\v z)\Big]_{\rho\rho\rho}&=-\ln\frac{\vee}{\Lambda}\int_{w,x,y,u}\frac{ig^3\,\rho^a(\v w)\,\rho^b(\v x)\,\rho^b(\v y)\,K(\v x,\v y,\v u)\,\Kv^i(\v w-\v z)}{8\sqrt2\,\pi^4\sqrt{k^+}}
\notag\\&=\mf A^{(1)a}_{2i}(k^+,-\v z)\,\mf N_2\,.
\end{align}
The operator ordering in Eq.~\eqref{eq:A3threerho}, with $\rho^a(\v w)$ to the left of the norm factor, is dictated by Eq.~\eqref{eq:lcpt}. It is important in what follows: since $\mf N_2\mf A^{(1)}\neq\mf A^{(1)}\mf N_2$, the opposite ordering would differ by $-[\mf N_2,\mf A^{(1)}]$, which is a two-$\rho$ term equal and opposite to the term $-2k^i\ln(\vee/\Lambda)$ of Eq.~\eqref{eq:A3tworho}.

The one-$\rho$, two-$\rho$ and three-$\rho$ parts of the wave function from which $\mf A^{(3)}_2$ is obtained agree with the NLO light-cone wave function of Lublinsky and Mulian~\cite{Lublinsky:2016meo} (eqs.~3.97--3.102 there), including the coefficient $11/3$, the constant $67/9-\pi^2/3$, the dependence on the longitudinal cut-offs and the three-$\rho$ term; the two-$\rho$ parts agree after the antisymmetrisation enforced by $f^{abc}\{\rho^b,\rho^c\}$. The coordinate-space forms~\eqref{eq:A3onerhoX}--\eqref{eq:A3threerho} go beyond the published material.

\paragraph{Two-gluon emission.} The coefficient $\mf B_2$ gives the amplitude for two soft gluons, with longitudinal momenta $k^+$ and $p^+$, at transverse positions $\v w$ and $\v z$:
\begin{align}\label{eq:B2}
\mf B^{bc}_{2ji}(k^+,-\v w;p^+,-\v z)=\ &\int_x\frac{ig^2f^{cba}\,\rho^a(\v x)\,\sqrt{p^+k^+}}{4\pi^2\big[p^+(\v x-\v z)^2+k^+(\v x-\v w)^2\big](p^++k^+)}
\notag\\&\times\Bigg[\frac{\delta^{ij}\big[(\v x-\v z)^2-(\v x-\v w)^2\big]}{2(\v z-\v w)^2}
\notag\\&\quad
+\frac{p^++k^+}{k^+}\Big(\frac{(x-z)^j(z-w)^i}{(\v z-\v w)^2}-\frac{(x-z)^j(x-w)^i}{2(\v x-\v w)^2}\Big)
\notag\\&\quad
+\frac{p^++k^+}{p^+}\Big(\frac{(x-w)^i(z-w)^j}{(\v z-\v w)^2}+\frac{(x-z)^j(x-w)^i}{2(\v x-\v z)^2}\Big)\Bigg]
\notag\\&\quad\times\Theta(\vee-k^+-p^+)
\notag\\&
-\int_{x,y}\frac{g^2\,\{\rho^c(\v y),\rho^b(\v x)\}\,\Kv^i(\v y-\v w)\,\Kv^j(\v x-\v z)}{8\pi^2\sqrt{p^+k^+}}
\notag\\&\quad\times\Theta(\vee-k^+)\,\Theta(\vee-p^+)\,,
\end{align}
with $k^+,p^+>\Lambda$. The first term, linear in $\rho$, arises from the emission of a gluon of momentum $k^++p^+$ which splits into two (triple-gluon vertex), from the instantaneous emission of the pair, and from the commutator part of the emission of two gluons by two charges (Figure~\ref{fig:vac_gg}); the parent gluon is a soft mode, hence $k^++p^+<\vee$. The second term is the independent emission of the two gluons by two valence charges; each gluon is separately a soft mode, and there is no condition on $k^++p^+$.
% CHOICE FOR THE AUTHORS: the support of the commutator part of the two-emission diagram (contained in the first term) is taken as k^+ + p^+ < vee, as for the splitting diagrams; by the kinematics of that diagram it could also be the full square Lambda<k^+,p^+<vee.

\paragraph{Scattering of a soft gluon on the valence charges.} The coefficient $\mf B_3$ describes the transition of a soft gluon of momentum $k^+$ at $\v z$ into a soft gluon of momentum $p^+$ at $\v u$ through its interaction with the valence charges: instantaneous interaction, emission of a second gluon followed by absorption of the first, and absorption of the gluon followed by re-emission (Figure~\ref{fig:g_g}). For $p^+>k^+$ the valence charges transfer longitudinal momentum to the gluon, for $p^+<k^+$ the gluon transfers momentum to the valence charges. With $\Theta_>\equiv\Theta(p^+-k^+-\Lambda)$ and $\Theta_<\equiv\Theta(k^+-\Lambda-p^+)$,
\begin{align}\label{eq:B3}
\mf B^{ba}_{3ji}(k^+,\v z;p^+,-\v u)=\ &\frac{ig^2f^{abc}}{(2\pi)^2}\int_x\frac{\sqrt{p^+k^+}\,(p^++k^+)}{(k^+-p^+)^2}\,\frac{\rho^c(\v x)\,\delta_{ij}}{p^+(\v x-\v u)^2-k^+(\v x-\v z)^2}
\notag\\&\qquad\times\big[\Theta_>+\Theta_<\big]
\notag\\&
+\frac{ig^2f^{abc}}{2\pi^2k^+}\int_x\frac{\sqrt{k^+p^+}}{k^+-p^+}
\notag\\&\qquad\times\frac{\big(k^+(x-z)^m-p^+(x-u)^m\big)\big((x-z)^l-(x-u)^l\big)}{k^+(\v x-\v z)^2-p^+(\v x-\v u)^2}\,\rho^c(\v x)
\notag\\&\qquad\times\Big(\delta_{il}\delta_{jm}-\frac{k^+}{p^+}\delta_{jl}\delta_{im}-\frac{k^+}{k^+-p^+}\delta_{ij}\delta_{lm}\Big)
\notag\\&\qquad\times\Bigg[-\frac{\Theta_>}{(\v u-\v z)^2}-\frac{p^+k^+\,\Theta_<}{\big(k^+(\v x-\v z)-p^+(\v x-\v u)\big)^2}\Bigg]
\notag\\&
-\frac{g^2}{4\pi^2\sqrt{p^+k^+}}\int_{x,y}\Kv^i(\v x-\v z)\,\Kv^j(\v y-\v u)\,\{\rho^a(\v x),\rho^b(\v y)\}
\notag\\&
-\frac{ig^2f^{abc}}{4\pi^2\sqrt{p^+k^+}}\int_x\Kv^i(\v x-\v z)\,\Kv^j(\v x-\v u)
\notag\\&\qquad\times\frac{k^+(\v x-\v z)^2+p^+(\v x-\v u)^2}{p^+(\v x-\v u)^2-k^+(\v x-\v z)^2}\,\rho^c(\v x)\,.
\end{align}
The double pole $1/(k^+-p^+)^2$ at the boundary $|p^+-k^+|=\Lambda$ cancels between the instantaneous and the non-instantaneous terms of the first two lines, so that no term proportional to $k^+/\Lambda$ arises. The vanishing denominators $p^+(\v x-\v u)^2=k^+(\v x-\v z)^2$ are principal values.

\paragraph{Gluon splitting.} The coefficient $\mf C_1$ is independent of $\rho$; it is the light-cone amplitude for the splitting of a soft gluon $(k^+,\v z)$ into two soft gluons of momenta $p^+$ and $k^+-p^+$ (Figure~\ref{fig:g_gg}):
\begin{align}\label{eq:C1}
\mf C^{bca}_{1jki}(k^+,\v z;p^+,\v x)=\ &-\frac{gf^{abc}}{2\pi\sqrt{2k^+}}\,\frac{\sqrt{p^+(k^+-p^+)}}{k^+}\,\delta^{(2)}\Big(\v z+\frac{p^+}{k^+}\v x\Big)\,\Kv^m(\v x)
\notag\\&\times\Big[\delta_{im}\delta_{jk}-\frac{k^+}{p^+}\delta_{jm}\delta_{ik}-\frac{k^+}{k^+-p^+}\delta_{km}\delta_{ij}\Big]\Theta(k^+-p^+-\Lambda)\,.
\end{align}
The $\delta$-function expresses the conservation of the transverse centre of mass $p^+\v x_1+(k^+-p^+)\v x_2=k^+\v x_0$ in the splitting.

\paragraph{Coefficients fixed by unitarity.} $\mf A^{(3)}_1$ and $\mf B_1$, which multiply annihilation operators, are not independent; they are given in terms of the coefficients above by Eqs.~\eqref{eq:A31unit} and~\eqref{eq:B1unit}. In the cross section they appear only through these relations.
